\subsection{A projector that cannot descend to the designated global form}
The determinant of one adapted basis does not by itself rule out
an alternative basis of the original tensor lattice. An actual
commutant projector gives a stronger test.

\begin{finding}[\fid{F49}: an intrinsic global-form obstruction]\label{f:globallattice}
Put
\[
 R=\Q[v,v^{-1}],\quad F=\Q(v),\quad
 \delta=v+v^{-1},\quad D=\delta^2,\quad h=D-1,
\]
and let
\[
 M_R=(R^2\otimes_R R^2)^{\otimes3}
\]
be the original tensor lattice with its actual nonstandard action.
There is a projector $p\in\End_F(M_F)$ which commutes with both
named right generators and all the named ordinary left operators,
is a sum of the diagonal matrix units in \fid{F43}, and is induced
by a literal coherent direct-summand projector in \fid{F46}, but
\[
 p(M_R)\not\subseteq M_R.
\]
Consequently no basis of this same lattice can make all the specified
\fid{F43} commutant matrix units integral. Algebraic coefficient
extensions with the prescribed integral closure do not remove this
obstruction.

Under the generic nonstandard Schur--Weyl identification, $p$ lies
in the Schur algebra but not in VIII's specified dual global form.
It is nevertheless compatible with the actual ordinary tensor bar
and preserves the two explicit monomial-coordinate lattices at
$v=0$ and $v=\infty$. No identification of those two lattices with
VIII's conjectural dual local crystal lattice is asserted.
\end{finding}

\begin{proof}
\par\smallskip\noindent\textbf{The actual ordinary projector.}\quad
Regroup $M_F$ as $V_F^{\otimes3}\otimes W_F^{\otimes3}$.
On the $V$-weight space with ordered monomials
$(m_{112},m_{121},m_{211})$, the ordinary upper canonical basis is
\[
 c_{112}=m_{112},\qquad
 c_{121}=m_{121}-v^{-1}m_{112},\qquad
 c_{211}=m_{211}-v^{-1}m_{121};
\]
see IV, Figure~1, p.~9, and Theorem~6.4, p.~34 \cite{IV}.
This change is unimodular over $R$.
The actual decomposition of \fid{F43} uses
\[
 b_1=v^2m_{112}+v m_{121}+m_{211},
\]
\[
 u_1=-(1+v^{-2})m_{112}+v m_{121}+m_{211},\qquad
 u_2=-v^{-1}m_{112}-v^{-2}m_{121}+\delta m_{211}.
\]
Their right-row coordinate matrix in the canonical basis is
\[
 \begin{pmatrix}
 h&\delta&1\\
 0&\delta&1\\
 0&1&\delta
 \end{pmatrix}.
\]
In particular $b_1=h c_{112}+u_1$.
The other three-dimensional weight space has the same rows,
with the last two interchanged; its first vector is
$b_2=v^2m_{122}+v m_{212}+m_{221}$.
Together with $b_0=m_{111}$ and $b_3=m_{222}$ these give
\[
 V_F^{\otimes3}
 =(W_3\otimes\epsilon)\oplus(W_1\otimes U).
\]
The actual coproduct calculation gives
\[
 Fb_0=b_1,\quad Fb_1=\delta b_2,\quad Fb_2=h b_3,\quad Fb_3=0,
\]
\[
 Eb_0=0,\quad Eb_1=h b_0,\quad Eb_2=\delta b_1,\quad Eb_3=b_2.
\]
The two $W_1$ copies have $Fg_+=g_-$ and $Eg_-=g_+$,
with the other actions zero. The ordinary right Hecke action
is scalar $\delta$ on $\epsilon$, and on $U$ is
\[
 c_1=\begin{pmatrix}\delta&0\\1&0\end{pmatrix},
 \qquad
 c_2=\begin{pmatrix}0&1\\0&\delta\end{pmatrix}.
\]
Thus these are actual invariant summands.

Let $p_V$ project onto $W_3\otimes\epsilon$ along $W_1\otimes U$.
Its two weight-space matrices, in canonical and monomial
coordinates respectively, are
\[
 [p_V]_{\mathrm{can}}=
 \begin{pmatrix}
 1&\delta/h&1/h\\
 0&0&0\\
 0&0&0
 \end{pmatrix},
 \qquad
 [p_V]_{\mathrm{mon}}=
 \frac1h
 \begin{pmatrix}
 v^2&v&1\\
 v&1&v^{-1}\\
 1&v^{-1}&v^{-2}
 \end{pmatrix}.
\]
The first follows from the preceding coordinate matrix.
For the second put $u=(v,1,v^{-1})$.
Its numerator is $u^{\mathsf t}u$, and $u u^{\mathsf t}=h$,
so the matrix is idempotent. It fixes $b_1=v u$ and kills
$u_1,u_2$. The lower weight space has the same matrix;
on each endpoint weight space $p_V$ is one.

In particular
\[
 p_V(m_{112})
 =c_{112}+\frac{\delta}{h}c_{121}+\frac1h c_{211}
 =\frac{v^2m_{112}+v m_{121}+m_{211}}h.
\]
Since $h=v^2+1+v^{-2}$ is a nonunit of $R$, the coefficient
$1/h$ is not in $R$. Hence $p_V$ does not preserve the
original lattice.

Now put $p=p_V\otimes1$ on $M_F$. It commutes with
\[
 P_i=c_i^V\otimes c_i^W+
       (c_i^V-\delta)\otimes(c_i^W-\delta)
\]
because $p_V$ commutes with $c_i^V$.
It also commutes with the ordinary left $E,F,K^{\pm1}$
operators on both factors. The coefficient of
$m_{211}\otimes m_{111}$ in
$p(m_{112}\otimes m_{111})$ is again $1/h$, proving
$p\notin\End_R(M_R)$.
If any $R$-basis made its matrix integral, it would preserve
$M_R$, a contradiction. This is independent of the determinant
of any displayed change of basis.

If $1/h$ were integral over $R$, multiply a monic integral
equation by the appropriate power of $h$ to obtain $1\in hR$.
This is impossible. Therefore $1/h$ does not lie in the
integral closure of $R$ in any algebraic extension of $F$.
The proof is unchanged after adjoining $i$.
If a parameter convention first adjoins $v$ as a square root,
integrality over the smaller Laurent ring would imply
integrality over $R$, so it cannot repair this coefficient.

\par\smallskip\noindent\textbf{The specified commutant and coherent projector.}\quad
The multiplicity decomposition in \fid{F43} is
\[
 \begin{aligned}
 M={}&(W_3\otimes W_3\oplus W_1\otimes W_1)\otimes\epsilon_+\\
 &\oplus(W_1\otimes W_1)\otimes\epsilon_-\\
 &\oplus(W_3\otimes W_1\oplus W_1\otimes W_3)\otimes N(\delta)\\
 &\oplus(W_1\otimes W_1)\otimes N(D-2).
 \end{aligned}
\]
The multiplicity ranks are $20,4,16,4$.
The operator $p$ selects the first sixteen multiplicities in
the $\epsilon_+$ block and the first eight in the $N(\delta)$
block, and kills the others. It is therefore a sum of diagonal
matrix units of the full commutant
$M_{20}\times M_4\times M_{16}\times M_4$.
If all those units preserved $M_R$, so would $p$.

On \fid{F46}'s coefficient base $S$, retain its framed object
\[
 K=S_{-1}^{20}\oplus S_i^4\oplus L_1^{16}\oplus L_2^4.
\]
Its direct-summand projector
\[
 g=\operatorname{diag}(I_{16},0_4)\oplus0_4
   \oplus\operatorname{diag}(I_8,0_8)\oplus0_4
\]
is an actual endomorphism of both the coefficient module and
its corresponding finite-support coherent sheaf on $\mathbf P^1_S$.
It commutes with the coordinate $T$ and the chosen Frobenius.
Under the proved postcomposition isomorphism
\[
 \End_F(K)\simeq
 \End_A\operatorname{Hom}_F(H,K)
\]
it induces $p$, with the multiplicity framing above.
The branch comparison changes the two-dimensional right factor
only, and hence leaves these multiplicity selectors unchanged.
The obstruction concerns descent of this particular actual
projector; the coherent realization itself has the projector.

\par\smallskip\noindent\textbf{Bar and the two boundary tests.}\quad
The canonical vectors above are fixed by the actual tensor bar.
All coefficients in their displayed adapted coordinate matrix
are bar invariant. The endpoints and lower weight analogue
have the same property. Therefore both summands defining $p_V$
are bar stable and $p_V$, hence $p$, commutes with the actual bar.
This does not require the tensor monomials to be bar fixed.

For the framed coherent comparison one may send $v$ to $v^{-1}$
and fix $i,x,T$ in \fid{F46}'s coordinate presentations.
The coefficients $\delta,D,h,\Delta,s,\lambda$ are invariant.
The comparison of a right block $N(a)$, $a=\delta,D-2$, with
the corresponding branch presentation has invariant matrix
\[
 S_a=\begin{pmatrix}1&0\\-(a/D)^2&a/D\end{pmatrix},
\]
which is invertible on that coefficient base.
With the selected multiplicity frames this gives the compatible
bar on the Hom realization, and $g$ is fixed.
It is a framed comparison, not a unique global-basis selection.

Let $A_0,A_\infty$ be the rings of rational functions regular
at zero and infinity. The nonzero entries of the monomial
projector are $v^j/h$ with $-2\leq j\leq2$.
Since
\[
 \frac{v^j}{h}=\frac{v^{j+2}}{v^4+v^2+1},
\]
their orders at zero and infinity are respectively $j+2$
and $2-j$, both nonnegative. Thus $p$ preserves both explicit
monomial-coordinate lattices $M_{A_0}$ and $M_{A_\infty}$.
The canonical-coordinate entries $1,\delta/h,1/h$ are also
regular at both boundaries.
This does not assert membership in VIII's conjectural dual
local crystal lattice. It proves that these two concrete
boundary tests and actual bar compatibility do not imply
the global Laurent condition: the latter fails at finite
nonzero zeros of $h$.

\par\smallskip\noindent\textbf{VIII's precise global form.}\quad
Write $\mathcal S$ for the nonstandard coordinate algebra.
VIII defines $\mathcal S_{\Q}$ as the $K_{\Q}$-subring generated
by the generic matrix entries $u^i_j$, with $K_{\Q}$ the
integral closure of the Laurent parameter ring.
Its degree-three Schur algebra is
$\mathcal A=(\mathcal S_3)^*$ and its specified dual global form is
\[
 \mathcal A_{\Q}
 =\{a\in\mathcal A:
       \langle a,\mathcal S_{3,\Q}\rangle\subseteq K_{\Q}\}.
\]
These are VIII, \S2.2, pp.~7, 9--10, and \S2.3, p.~11
\cite{VIII}.
The global-form part of its Proposition~2.4, pp.~12--13,
has the following direct proof in this situation.
Write the tensor coaction in the original basis as
\[
 e_J\longmapsto\sum_I e_I\otimes u_{I,J}.
\]
Each $u_{I,J}$ is a degree-three product of generating
coordinate entries and belongs to $\mathcal S_{3,\Q}$.
The matrix entries of the induced operator $a$ are
$\langle a,u_{I,J}\rangle$, hence belong to $K_{\Q}$ whenever
$a\in\mathcal A_{\Q}$. Consequently
\[
 \eta(\mathcal A_{\Q})
 \subseteq\End_{K_{\Q}}(X_{K_{\Q}}^{\otimes3}).
\]
For $X=V\otimes W$ under the direct-product ordinary quantum
group, its chosen upper canonical basis is the four elementary
products. This is the same original tensor form after scalar
extension.

Generic nonstandard Schur--Weyl duality identifies the full
commutant with $\eta(\mathcal A)$, so $p$ belongs to that algebra.
The remark immediately after IV Theorem~12.1 corrects its
tensor-dual convention relative to \S\S8--10.
Transposing that identification preserves the global entry
criterion; our monomial projector matrices are symmetric,
so the present $p$ is unchanged by that transpose.
Its coefficient $1/h$ and the preceding integral-closure
argument prove
\[
 p\notin\eta(\mathcal A_{\Q}).
\]
The localized matrix-unit basis therefore cannot, under this
actual action, be identified with a basis of the designated
global form; not all its elements even belong to that form.
The global inclusion just used has been proved directly.
The printed proof of VIII Proposition~2.4 is a placeholder;
its additional local-lattice inclusion is conditional on
Conjecture~2.2 and has not been assumed here.
\end{proof}

\paragraph{The exact boundary of this obstruction.}
IV already constructs a based rank-three example in
Example~19.2, pp.~132--133. Its paragraph after Conjecture~19.1,
p.~131, explicitly distinguishes the desired integral form and
local lattice from the original tensor forms \cite{IV}.
Our obstruction is compatible with that construction.
It concerns simultaneous retention of the original global tensor
lattice and the full specified \fid{F43}/\fid{F46} commutant
projectors. It does not obstruct an alternate basis with only
the named nonstandard generators, a changed integral form, a
smaller integral commutant, or VIII's conjectured minimum-degree
globalization within its designated form.

% Separate paragraph for Part III.
\paragraph{\fid{F49} in Grothendieck's language.}
The original lattice is a vector bundle over $\Spec R$.
The rational idempotent $p$ extends and splits that bundle
over the open set $h\ne0$, but its explicit $1/h$ coefficient
prevents extension as an endomorphism of this bundle over
the whole parameter scheme. The precise affine dictionary is
EGA~I, Theorem~1.3.7, p.~86, Corollary~1.3.8, p.~87, and
Theorem~1.4.1, pp.~90--92.
The perfect direct-summand framework is SGA~6, I,
Definition~4.7, pp.~121--122, and Example~4.8, p.~122;
cote~90, scan~17 is an audited equivariant-module anchor.
None of those sources is claimed to contain this projector
calculation or the comparison with VIII.
The coherent model retains the actual framed bar and projector;
descent of its full specified localized projector system to
the designated global form fails under the actual tensor
comparison. Changing the integral form or restricting the
operator algebra remains possible. Purity of the separately
chosen Weil coefficients does not supply that comparison.
